\section{Formula Kit --- Emit}
%%% generated by the blueprint pipeline from TCSlib.Complexity.FormulaKit.Emit
\subsection*{Formula Kit --- Emit: Introduction}

Declarations of the Lean module \texttt{TCSlib.Complexity.FormulaKit.Emit}.

\subsection*{Formula Kit --- Emit: Results}

\begin{definition}[Clause-list encoding fragment]
\lean{Complexity.FormulaKit.fragment}
\label{Complexity.FormulaKit.fragment}
\leanok
\uses{Std.Sat.CNF.serializeClause}
For a CNF formula $\varphi$ over natural-number variables, its fragment is the Boolean
list obtained by replacing each clause $C$, in order, with the list consisting of `true`
followed by the serialization of $C$, and concatenating these lists.
\end{definition}

\begin{theorem}[Formula encoding and clause fragment]
\lean{Complexity.FormulaKit.serialize_eq_fragment}
\label{Complexity.FormulaKit.serialize_eq_fragment}
\leanok
\uses{Complexity.FormulaKit.fragment, Std.Sat.CNF.serialize}
\difficulty{1}
Let $\varphi=(C_1,\ldots,C_m)$ be a CNF formula over variables indexed by the natural
numbers. Form a bit string by concatenating, for each clause $C_i$, a true bit followed
by its clause encoding; then the encoding of $\varphi$ is exactly this string followed
by one false bit.
\end{theorem}

\begin{theorem}[Fragment of a concatenation]
\lean{Complexity.FormulaKit.fragment_append}
\label{Complexity.FormulaKit.fragment_append}
\leanok
\uses{Complexity.FormulaKit.fragment}
\difficulty{1}
For any two CNF formulas $\varphi$ and $\psi$ over natural-number variables, the binary
fragment of their concatenation is the concatenation of their binary fragments.
\end{theorem}

\begin{theorem}[Fragment of a flattened list]
\lean{Complexity.FormulaKit.fragment_flatMap}
\label{Complexity.FormulaKit.fragment_flatMap}
\leanok
\uses{Complexity.FormulaKit.fragment, Complexity.FormulaKit.fragment_append}
\difficulty{4}
For any type $\beta$, any finite list $xs$ of elements of $\beta$, and any map $G$
assigning a CNF formula over natural-number variables to each element of $\beta$, the
binary fragment of the formula obtained by concatenating the formulas $G(b)$ for $b$ in
$xs$ is the concatenation, in list order, of the binary fragments of those formulas.
\end{theorem}

\begin{theorem}[Polynomial-Time Equivalence of Formula Encodings]
\lean{Complexity.FormulaKit.polyTimeComputable_fragment_iff_serialize}
\label{Complexity.FormulaKit.polyTimeComputable_fragment_iff_serialize}
\leanok
\uses{Complexity.FormulaKit.fragment, Complexity.PolyTimeComputable, Std.Sat.CNF.serialize}
For every function $F$ from bit strings to CNF formulas over natural-number-indexed
variables, the function sending $x$ to the fragment encoding of $F(x)$ is
polynomial-time computable if and only if the function sending $x$ to the binary
serialization of $F(x)$ is polynomial-time computable.
\end{theorem}

\begin{theorem}[Polynomial-Time Indexed Concatenation]
\lean{Complexity.FormulaKit.polyTimeComputable_flatMap_range}
\label{Complexity.FormulaKit.polyTimeComputable_flatMap_range}
\leanok
\uses{Complexity.PolyTimeComputable, Complexity.pairFstD, Complexity.pairSndD}
Let $c$ assign a natural number to each bit string, and let $g(x,i)$ be a bit string for
each bit string $x$ and natural number $i$. If the function sending $x$ to the unary
string of $c(x)$ ones is polynomial-time computable, and the function sending a pair
encoding of bit strings $a,b$ to $g(a,|b|)$ is polynomial-time computable, then the
function sending $x$ to the concatenation of $g(x,i)$ over all integers $i$ with $0 \le
i < c(x)$ is polynomial-time computable.
\end{theorem}

\begin{theorem}[Polynomial-Time Serialization of Indexed Formula Concatenation]
\lean{Complexity.FormulaKit.polyTimeComputable_serialize_flatMap}
\label{Complexity.FormulaKit.polyTimeComputable_serialize_flatMap}
\leanok
\uses{Complexity.FormulaKit.fragment, Complexity.FormulaKit.fragment_flatMap, Complexity.FormulaKit.polyTimeComputable_flatMap_range, Complexity.FormulaKit.serialize_eq_fragment, Complexity.PolyTimeComputable, Complexity.PolyTimeComputable.append, Complexity.pairFstD, Complexity.pairSndD, Complexity.polyTimeComputable_const, Std.Sat.CNF.serialize}
\difficulty{4}
Let $c$ assign a natural number to each bit string, and let $G(x,i)$ be a CNF formula
over natural-number-indexed variables for each bit string $x$ and natural number $i$. If
the function sending $x$ to the unary string of $c(x)$ ones is polynomial-time
computable, and the function sending a pair encoding of bit strings $a,b$ to the
fragment encoding of $G(a,|b|)$ is polynomial-time computable, then the function sending
$x$ to the binary serialization of the CNF formula formed by concatenating $G(x,i)$ for
all $i$ with $0\le i<c(x)$ is polynomial-time computable.
\end{theorem}

\begin{theorem}[Polynomial-Time Encoding After Variable Relabeling]
\lean{Complexity.FormulaKit.polyTimeComputable_fragment_relabel}
\label{Complexity.FormulaKit.polyTimeComputable_fragment_relabel}
\leanok
\uses{Complexity.FormulaKit.fragment, Complexity.PolyTimeComputable, Std.Sat.CNF.numVars}
Let $\varphi$ be a CNF formula over natural-number-indexed variables, and let $w(z,j)$
assign a natural number to each bit string $z$ and natural number $j$. Suppose that for
every $j<N(\varphi)$, where $N(\varphi)$ is the variable count of $\varphi$, the unary
encoding of $w(z,j)$ is polynomial-time computable. Then the function sending $z$ to the
fragment encoding of the formula obtained by relabeling each variable index $j$ of
$\varphi$ by $w(z,j)$ is polynomial-time computable.
\end{theorem}

\begin{theorem}[Polynomial-Time Encoding of a Gate Fragment]
\lean{Complexity.FormulaKit.polyTimeComputable_fragment_gate}
\label{Complexity.FormulaKit.polyTimeComputable_fragment_gate}
\leanok
\uses{Complexity.FormulaKit.fragment, Complexity.FormulaKit.gate, Complexity.FormulaKit.polyTimeComputable_fragment_relabel, Complexity.FormulaKit.template, Complexity.FormulaKit.template_numVars_le, Complexity.PolyTimeComputable}
\difficulty{2}
Let $\ell$ be a natural number, let $p$ be a Boolean function on bit strings of length
$\ell$, and let $w(z,j)$ assign a natural number to each bit string $z$ and natural
number $j$. Suppose that for every $j<\ell$, the function sending $z$ to the unary
string of $w(z,j)$ ones is polynomial-time computable. Then the function that sends $z$
to the binary fragment encoding of the CNF gate obtained by relabeling the truth-table
clauses for $p$ using the wire indices $w(z,j)$ is polynomial-time computable.
\end{theorem}

\begin{theorem}[Polynomial-Time Encoding of a Selected Gate Fragment]
\lean{Complexity.FormulaKit.polyTimeComputable_fragment_gate_select}
\label{Complexity.FormulaKit.polyTimeComputable_fragment_gate_select}
\leanok
\uses{Complexity.FormulaKit.fragment, Complexity.FormulaKit.gate, Complexity.PolyTimeComputable}
Let $\kappa$ be a finite set with decidable equality, let $\ell$ be a natural number,
and let $p(\kappa_0,u)$ be a Boolean function of a type $\kappa_0\in\kappa$ and a bit
string $u$ of length $\ell$. Let $k(z)\in\kappa$ be a function of bit strings, and let
$w(z,j)$ assign a natural number to each bit string and natural number $j$. Suppose that
for each $\kappa_0\in\kappa$, the function returning the one-bit indicator of
$k(z)=\kappa_0$ is polynomial-time computable, and that for every $j<\ell$, the unary
encoding of $w(z,j)$ is polynomial-time computable. Then the function sending $z$ to the
binary fragment encoding of the CNF gate for $p(k(z),\cdot)$, with wire indices
$w(z,j)$, is polynomial-time computable.
\end{theorem}

\begin{theorem}[Polynomial-Time Unary Addition]
\lean{Complexity.FormulaKit.polyTimeComputable_unary_add}
\label{Complexity.FormulaKit.polyTimeComputable_unary_add}
\leanok
\uses{Complexity.PolyTimeComputable, Complexity.PolyTimeComputable.append}
\difficulty{1}
Let $f$ and $g$ assign natural numbers to bit strings. If the functions sending $x$ to
the unary strings of $f(x)$ ones and of $g(x)$ ones are polynomial-time computable, then
the function sending $x$ to the unary string of $f(x)+g(x)$ ones is polynomial-time
computable.
\end{theorem}

\begin{theorem}[Polynomial-time unary truncated subtraction]
\lean{Complexity.FormulaKit.polyTimeComputable_unary_sub}
\label{Complexity.FormulaKit.polyTimeComputable_unary_sub}
\leanok
\uses{Complexity.PolyTimeComputable}
Let $f,g$ be functions from finite bit strings to the natural numbers. If the unary
string functions $x\mapsto 1^{f(x)}$ and $x\mapsto 1^{g(x)}$ are polynomial-time
computable, then so is $x\mapsto 1^{\max(f(x)-g(x),0)}$, with subtraction interpreted as
truncated natural-number subtraction.
\end{theorem}

\begin{theorem}[Polynomial-time unary multiplication]
\lean{Complexity.FormulaKit.polyTimeComputable_unary_mul}
\label{Complexity.FormulaKit.polyTimeComputable_unary_mul}
\leanok
\uses{Complexity.PolyTimeComputable}
Let $f,g$ be functions from finite bit strings to the natural numbers. If the unary
string functions $x\mapsto 1^{f(x)}$ and $x\mapsto 1^{g(x)}$ are polynomial-time
computable, then so is $x\mapsto 1^{f(x)g(x)}$.
\end{theorem}

\begin{theorem}[Polynomial-time unary powers]
\lean{Complexity.FormulaKit.polyTimeComputable_unary_pow}
\label{Complexity.FormulaKit.polyTimeComputable_unary_pow}
\leanok
\uses{Complexity.FormulaKit.polyTimeComputable_unary_mul, Complexity.PolyTimeComputable, Complexity.polyTimeComputable_const}
\difficulty{4}
Let $f$ be a function from finite bit strings to the natural numbers, and let
$k\in\mathbb N$. If the unary string function $x\mapsto 1^{f(x)}$ is polynomial-time
computable, then so is $x\mapsto 1^{f(x)^k}$.
\end{theorem}

\begin{theorem}[Polynomial-time unary comparison]
\lean{Complexity.FormulaKit.polyTimeComputable_unary_le}
\label{Complexity.FormulaKit.polyTimeComputable_unary_le}
\leanok
\uses{Complexity.PolyTimeComputable}
Let $f,g$ be functions from finite bit strings to the natural numbers. If the unary
string functions $x\mapsto 1^{f(x)}$ and $x\mapsto 1^{g(x)}$ are polynomial-time
computable, then so is the one-bit string function that returns true exactly when
$f(x)\le g(x)$.
\end{theorem}

\begin{theorem}[Polynomial-time unary division]
\lean{Complexity.FormulaKit.polyTimeComputable_unary_div}
\label{Complexity.FormulaKit.polyTimeComputable_unary_div}
\leanok
\uses{Complexity.PolyTimeComputable}
Let $f,g$ be functions from finite bit strings to the natural numbers. If the unary
string functions $x\mapsto 1^{f(x)}$ and $x\mapsto 1^{g(x)}$ are polynomial-time
computable, then so is $x\mapsto 1^{\lfloor f(x)/g(x)\rfloor}$, where $1^m$ denotes a
string of $m$ true bits and natural-number division gives quotient zero when the divisor
is zero.
\end{theorem}

\begin{theorem}[Polynomial-time unary remainder]
\lean{Complexity.FormulaKit.polyTimeComputable_unary_mod}
\label{Complexity.FormulaKit.polyTimeComputable_unary_mod}
\leanok
\uses{Complexity.PolyTimeComputable}
Let $f,g$ be functions from finite bit strings to the natural numbers. If the unary
string functions $x\mapsto 1^{f(x)}$ and $x\mapsto 1^{g(x)}$ are polynomial-time
computable, then so is $x\mapsto 1^{f(x)\bmod g(x)}$, where the remainder is taken in
natural-number arithmetic (and is $f(x)$ when $g(x)=0$).
\end{theorem}

\begin{theorem}[Polynomial-Time Serialization After Additive Relabeling]
\lean{Complexity.FormulaKit.polyTimeComputable_serialize_relabel_add}
\label{Complexity.FormulaKit.polyTimeComputable_serialize_relabel_add}
\leanok
\uses{Complexity.PolyTimeComputable, Std.Sat.CNF.serialize}
Let $F$ map bit strings to CNF formulas over natural-number-indexed variables, and let
$o$ map bit strings to natural numbers. If the function sending $x$ to the binary
serialization of $F(x)$ is polynomial-time computable, and the function sending $x$ to
the unary string of $o(x)$ ones is polynomial-time computable, then the function sending
$x$ to the binary serialization of the formula obtained from $F(x)$ by replacing each
variable index $j$ by $o(x)+j$ is polynomial-time computable.
\end{theorem}
